% =====================================================================
%  tikz_common_styles.tex
%  Stili condivisi per i diagrammi di architettura (Cap. 4 — Metodo).
%
%  Caricato da main.tex con % =====================================================================
%  tikz_common_styles.tex
%  Stili condivisi per i diagrammi di architettura (Cap. 4 — Metodo).
%
%  Caricato da main.tex con % =====================================================================
%  tikz_common_styles.tex
%  Stili condivisi per i diagrammi di architettura (Cap. 4 — Metodo).
%
%  Caricato da main.tex con % =====================================================================
%  tikz_common_styles.tex
%  Stili condivisi per i diagrammi di architettura (Cap. 4 — Metodo).
%
%  Caricato da main.tex con \input{tikz_common_styles} DOPO
%  \usepackage{tikz}. Le librerie sono ridichiarate qui per rendere il
%  file autonomo: \usetikzlibrary e' idempotente, non rompe nulla se
%  le stesse librerie sono gia' caricate in main.tex.
%
%  CRITERI TIPOGRAFICI
%  - Palette desaturata: scala di grigi + un solo blu-grigio a tre
%    luminanze per distinguere i livelli. Nessun colore saturo.
%  - Spigoli vivi (sharp corners): convenzione dei diagrammi a blocchi
%    in letteratura ingegneristica, meno "slide" dei bordi arrotondati.
%  - Larghezze dei box selezionabili con chiavi dedicate (wnarrow /
%    wmed / wwide) invece di essere riscritte a mano nelle figure:
%    il controllo del text wrapping resta centralizzato in questo file.
%  - Etichette degli archi su fondo bianco (fill=white) per evitare la
%    sovrapposizione con la linea e con i bordi dei box adiacenti.
% =====================================================================

\usetikzlibrary{positioning, arrows.meta, fit, backgrounds, calc}

% ---------------------------------------------------------------------
% Palette
% ---------------------------------------------------------------------
\definecolor{lawsInk}{gray}{0.22}          % bordi primari, testo etichette
\definecolor{lawsInkSoft}{gray}{0.45}      % archi, bordi secondari
\definecolor{lawsL1}{RGB}{231,237,243}     % Livello 1 — Visione
\definecolor{lawsL2}{RGB}{240,242,245}     % Livello 2 — Simulazione
\definecolor{lawsL3}{RGB}{247,248,249}     % Livello 3 — Metadati
\definecolor{lawsNeutral}{RGB}{247,247,247}% blocchi di processo
\definecolor{lawsAccent}{RGB}{224,230,237} % nodo decisionale (piu' scuro)

% ---------------------------------------------------------------------
% Titolo interno ai box: una sola definizione, spaziatura uniforme.
% Usare SEMPRE come primo elemento del testo del nodo, mai come ultimo
% (il \\ finale su nodo vuoto genera errore).
% ---------------------------------------------------------------------
\newcommand{\lawstitle}[1]{{\bfseries #1}\\[2.5pt]}

% ---------------------------------------------------------------------
% Stili
% ---------------------------------------------------------------------
\tikzset{%
  % --- classe base: qui si calibra il wrapping, non nelle figure ---
  lawsbox/.style={
    draw=lawsInk,
    line width=0.5pt,
    rectangle,
    sharp corners,
    fill=white,
    text width=46mm,          % default; sovrascritto da wnarrow/wmed/wwide
    align=center,
    inner xsep=6pt,           % respiro orizzontale: evita che il grassetto
    inner ysep=5.5pt,         % del titolo tocchi il bordo
    minimum height=13mm,      % pavimento, non vincolo: i box crescono
    font=\footnotesize,
    text=black,
  },
  % --- larghezze selezionabili ---
  wnarrow/.style={text width=40mm},
  wmed/.style={text width=48mm},
  wwide/.style={text width=82mm},
%
  % --- livelli architetturali ---
  vision/.style={lawsbox, fill=lawsL1},
  sim/.style={lawsbox, fill=lawsL2},
  meta/.style={lawsbox, fill=lawsL3},
%
  % --- blocchi di processo (pipeline di ottimizzazione) ---
  optim/.style={lawsbox, fill=lawsNeutral},
%
  % --- interfaccia di interscambio: bordo tratteggiato = artefatto su
  %     disco, non componente di calcolo ---
  bridge/.style={lawsbox, fill=white, draw=lawsInk, dashed, line width=0.6pt},
%
  % --- canali della fusione ---
  osint/.style={lawsbox, fill=lawsL2},
  behav/.style={lawsbox, fill=lawsL3},
%
  % --- nodi terminali: bordo piu' spesso, non colore diverso ---
  fusion/.style={lawsbox, fill=lawsNeutral, line width=0.9pt},
  decision/.style={lawsbox, fill=lawsAccent, line width=0.9pt},
%
  % --- archi ---
  arr/.style={
    -{Stealth[length=2mm, width=1.4mm]},
    draw=lawsInkSoft,
    line width=0.6pt,
  },
  busline/.style={draw=lawsInkSoft, line width=0.6pt},
%
  % --- etichette degli archi ---
  arrlabel/.style={
    font=\scriptsize,
    text=lawsInk,
    inner sep=1.5pt,
    fill=white,
  },
  % --- annotazione laterale (statuto epistemico, note di lettura) ---
  sidenote/.style={
    font=\scriptsize\itshape,
    text=lawsInkSoft,
    align=left,
    inner sep=1pt,
  },
}
% DOPO
%  \usepackage{tikz}. Le librerie sono ridichiarate qui per rendere il
%  file autonomo: \usetikzlibrary e' idempotente, non rompe nulla se
%  le stesse librerie sono gia' caricate in main.tex.
%
%  CRITERI TIPOGRAFICI
%  - Palette desaturata: scala di grigi + un solo blu-grigio a tre
%    luminanze per distinguere i livelli. Nessun colore saturo.
%  - Spigoli vivi (sharp corners): convenzione dei diagrammi a blocchi
%    in letteratura ingegneristica, meno "slide" dei bordi arrotondati.
%  - Larghezze dei box selezionabili con chiavi dedicate (wnarrow /
%    wmed / wwide) invece di essere riscritte a mano nelle figure:
%    il controllo del text wrapping resta centralizzato in questo file.
%  - Etichette degli archi su fondo bianco (fill=white) per evitare la
%    sovrapposizione con la linea e con i bordi dei box adiacenti.
% =====================================================================

\usetikzlibrary{positioning, arrows.meta, fit, backgrounds, calc}

% ---------------------------------------------------------------------
% Palette
% ---------------------------------------------------------------------
\definecolor{lawsInk}{gray}{0.22}          % bordi primari, testo etichette
\definecolor{lawsInkSoft}{gray}{0.45}      % archi, bordi secondari
\definecolor{lawsL1}{RGB}{231,237,243}     % Livello 1 — Visione
\definecolor{lawsL2}{RGB}{240,242,245}     % Livello 2 — Simulazione
\definecolor{lawsL3}{RGB}{247,248,249}     % Livello 3 — Metadati
\definecolor{lawsNeutral}{RGB}{247,247,247}% blocchi di processo
\definecolor{lawsAccent}{RGB}{224,230,237} % nodo decisionale (piu' scuro)

% ---------------------------------------------------------------------
% Titolo interno ai box: una sola definizione, spaziatura uniforme.
% Usare SEMPRE come primo elemento del testo del nodo, mai come ultimo
% (il \\ finale su nodo vuoto genera errore).
% ---------------------------------------------------------------------
\newcommand{\lawstitle}[1]{{\bfseries #1}\\[2.5pt]}

% ---------------------------------------------------------------------
% Stili
% ---------------------------------------------------------------------
\tikzset{%
  % --- classe base: qui si calibra il wrapping, non nelle figure ---
  lawsbox/.style={
    draw=lawsInk,
    line width=0.5pt,
    rectangle,
    sharp corners,
    fill=white,
    text width=46mm,          % default; sovrascritto da wnarrow/wmed/wwide
    align=center,
    inner xsep=6pt,           % respiro orizzontale: evita che il grassetto
    inner ysep=5.5pt,         % del titolo tocchi il bordo
    minimum height=13mm,      % pavimento, non vincolo: i box crescono
    font=\footnotesize,
    text=black,
  },
  % --- larghezze selezionabili ---
  wnarrow/.style={text width=40mm},
  wmed/.style={text width=48mm},
  wwide/.style={text width=82mm},
%
  % --- livelli architetturali ---
  vision/.style={lawsbox, fill=lawsL1},
  sim/.style={lawsbox, fill=lawsL2},
  meta/.style={lawsbox, fill=lawsL3},
%
  % --- blocchi di processo (pipeline di ottimizzazione) ---
  optim/.style={lawsbox, fill=lawsNeutral},
%
  % --- interfaccia di interscambio: bordo tratteggiato = artefatto su
  %     disco, non componente di calcolo ---
  bridge/.style={lawsbox, fill=white, draw=lawsInk, dashed, line width=0.6pt},
%
  % --- canali della fusione ---
  osint/.style={lawsbox, fill=lawsL2},
  behav/.style={lawsbox, fill=lawsL3},
%
  % --- nodi terminali: bordo piu' spesso, non colore diverso ---
  fusion/.style={lawsbox, fill=lawsNeutral, line width=0.9pt},
  decision/.style={lawsbox, fill=lawsAccent, line width=0.9pt},
%
  % --- archi ---
  arr/.style={
    -{Stealth[length=2mm, width=1.4mm]},
    draw=lawsInkSoft,
    line width=0.6pt,
  },
  busline/.style={draw=lawsInkSoft, line width=0.6pt},
%
  % --- etichette degli archi ---
  arrlabel/.style={
    font=\scriptsize,
    text=lawsInk,
    inner sep=1.5pt,
    fill=white,
  },
  % --- annotazione laterale (statuto epistemico, note di lettura) ---
  sidenote/.style={
    font=\scriptsize\itshape,
    text=lawsInkSoft,
    align=left,
    inner sep=1pt,
  },
}
% DOPO
%  \usepackage{tikz}. Le librerie sono ridichiarate qui per rendere il
%  file autonomo: \usetikzlibrary e' idempotente, non rompe nulla se
%  le stesse librerie sono gia' caricate in main.tex.
%
%  CRITERI TIPOGRAFICI
%  - Palette desaturata: scala di grigi + un solo blu-grigio a tre
%    luminanze per distinguere i livelli. Nessun colore saturo.
%  - Spigoli vivi (sharp corners): convenzione dei diagrammi a blocchi
%    in letteratura ingegneristica, meno "slide" dei bordi arrotondati.
%  - Larghezze dei box selezionabili con chiavi dedicate (wnarrow /
%    wmed / wwide) invece di essere riscritte a mano nelle figure:
%    il controllo del text wrapping resta centralizzato in questo file.
%  - Etichette degli archi su fondo bianco (fill=white) per evitare la
%    sovrapposizione con la linea e con i bordi dei box adiacenti.
% =====================================================================

\usetikzlibrary{positioning, arrows.meta, fit, backgrounds, calc}

% ---------------------------------------------------------------------
% Palette
% ---------------------------------------------------------------------
\definecolor{lawsInk}{gray}{0.22}          % bordi primari, testo etichette
\definecolor{lawsInkSoft}{gray}{0.45}      % archi, bordi secondari
\definecolor{lawsL1}{RGB}{231,237,243}     % Livello 1 — Visione
\definecolor{lawsL2}{RGB}{240,242,245}     % Livello 2 — Simulazione
\definecolor{lawsL3}{RGB}{247,248,249}     % Livello 3 — Metadati
\definecolor{lawsNeutral}{RGB}{247,247,247}% blocchi di processo
\definecolor{lawsAccent}{RGB}{224,230,237} % nodo decisionale (piu' scuro)

% ---------------------------------------------------------------------
% Titolo interno ai box: una sola definizione, spaziatura uniforme.
% Usare SEMPRE come primo elemento del testo del nodo, mai come ultimo
% (il \\ finale su nodo vuoto genera errore).
% ---------------------------------------------------------------------
\newcommand{\lawstitle}[1]{{\bfseries #1}\\[2.5pt]}

% ---------------------------------------------------------------------
% Stili
% ---------------------------------------------------------------------
\tikzset{%
  % --- classe base: qui si calibra il wrapping, non nelle figure ---
  lawsbox/.style={
    draw=lawsInk,
    line width=0.5pt,
    rectangle,
    sharp corners,
    fill=white,
    text width=46mm,          % default; sovrascritto da wnarrow/wmed/wwide
    align=center,
    inner xsep=6pt,           % respiro orizzontale: evita che il grassetto
    inner ysep=5.5pt,         % del titolo tocchi il bordo
    minimum height=13mm,      % pavimento, non vincolo: i box crescono
    font=\footnotesize,
    text=black,
  },
  % --- larghezze selezionabili ---
  wnarrow/.style={text width=40mm},
  wmed/.style={text width=48mm},
  wwide/.style={text width=82mm},
%
  % --- livelli architetturali ---
  vision/.style={lawsbox, fill=lawsL1},
  sim/.style={lawsbox, fill=lawsL2},
  meta/.style={lawsbox, fill=lawsL3},
%
  % --- blocchi di processo (pipeline di ottimizzazione) ---
  optim/.style={lawsbox, fill=lawsNeutral},
%
  % --- interfaccia di interscambio: bordo tratteggiato = artefatto su
  %     disco, non componente di calcolo ---
  bridge/.style={lawsbox, fill=white, draw=lawsInk, dashed, line width=0.6pt},
%
  % --- canali della fusione ---
  osint/.style={lawsbox, fill=lawsL2},
  behav/.style={lawsbox, fill=lawsL3},
%
  % --- nodi terminali: bordo piu' spesso, non colore diverso ---
  fusion/.style={lawsbox, fill=lawsNeutral, line width=0.9pt},
  decision/.style={lawsbox, fill=lawsAccent, line width=0.9pt},
%
  % --- archi ---
  arr/.style={
    -{Stealth[length=2mm, width=1.4mm]},
    draw=lawsInkSoft,
    line width=0.6pt,
  },
  busline/.style={draw=lawsInkSoft, line width=0.6pt},
%
  % --- etichette degli archi ---
  arrlabel/.style={
    font=\scriptsize,
    text=lawsInk,
    inner sep=1.5pt,
    fill=white,
  },
  % --- annotazione laterale (statuto epistemico, note di lettura) ---
  sidenote/.style={
    font=\scriptsize\itshape,
    text=lawsInkSoft,
    align=left,
    inner sep=1pt,
  },
}
% DOPO
%  \usepackage{tikz}. Le librerie sono ridichiarate qui per rendere il
%  file autonomo: \usetikzlibrary e' idempotente, non rompe nulla se
%  le stesse librerie sono gia' caricate in main.tex.
%
%  CRITERI TIPOGRAFICI
%  - Palette desaturata: scala di grigi + un solo blu-grigio a tre
%    luminanze per distinguere i livelli. Nessun colore saturo.
%  - Spigoli vivi (sharp corners): convenzione dei diagrammi a blocchi
%    in letteratura ingegneristica, meno "slide" dei bordi arrotondati.
%  - Larghezze dei box selezionabili con chiavi dedicate (wnarrow /
%    wmed / wwide) invece di essere riscritte a mano nelle figure:
%    il controllo del text wrapping resta centralizzato in questo file.
%  - Etichette degli archi su fondo bianco (fill=white) per evitare la
%    sovrapposizione con la linea e con i bordi dei box adiacenti.
% =====================================================================

\usetikzlibrary{positioning, arrows.meta, fit, backgrounds, calc}

% ---------------------------------------------------------------------
% Palette
% ---------------------------------------------------------------------
\definecolor{lawsInk}{gray}{0.22}          % bordi primari, testo etichette
\definecolor{lawsInkSoft}{gray}{0.45}      % archi, bordi secondari
\definecolor{lawsL1}{RGB}{231,237,243}     % Livello 1 — Visione
\definecolor{lawsL2}{RGB}{240,242,245}     % Livello 2 — Simulazione
\definecolor{lawsL3}{RGB}{247,248,249}     % Livello 3 — Metadati
\definecolor{lawsNeutral}{RGB}{247,247,247}% blocchi di processo
\definecolor{lawsAccent}{RGB}{224,230,237} % nodo decisionale (piu' scuro)

% ---------------------------------------------------------------------
% Titolo interno ai box: una sola definizione, spaziatura uniforme.
% Usare SEMPRE come primo elemento del testo del nodo, mai come ultimo
% (il \\ finale su nodo vuoto genera errore).
% ---------------------------------------------------------------------
\newcommand{\lawstitle}[1]{{\bfseries #1}\\[2.5pt]}

% ---------------------------------------------------------------------
% Stili
% ---------------------------------------------------------------------
\tikzset{%
  % --- classe base: qui si calibra il wrapping, non nelle figure ---
  lawsbox/.style={
    draw=lawsInk,
    line width=0.5pt,
    rectangle,
    sharp corners,
    fill=white,
    text width=46mm,          % default; sovrascritto da wnarrow/wmed/wwide
    align=center,
    inner xsep=6pt,           % respiro orizzontale: evita che il grassetto
    inner ysep=5.5pt,         % del titolo tocchi il bordo
    minimum height=13mm,      % pavimento, non vincolo: i box crescono
    font=\footnotesize,
    text=black,
  },
  % --- larghezze selezionabili ---
  wnarrow/.style={text width=40mm},
  wmed/.style={text width=48mm},
  wwide/.style={text width=82mm},
%
  % --- livelli architetturali ---
  vision/.style={lawsbox, fill=lawsL1},
  sim/.style={lawsbox, fill=lawsL2},
  meta/.style={lawsbox, fill=lawsL3},
%
  % --- blocchi di processo (pipeline di ottimizzazione) ---
  optim/.style={lawsbox, fill=lawsNeutral},
%
  % --- interfaccia di interscambio: bordo tratteggiato = artefatto su
  %     disco, non componente di calcolo ---
  bridge/.style={lawsbox, fill=white, draw=lawsInk, dashed, line width=0.6pt},
%
  % --- canali della fusione ---
  osint/.style={lawsbox, fill=lawsL2},
  behav/.style={lawsbox, fill=lawsL3},
%
  % --- nodi terminali: bordo piu' spesso, non colore diverso ---
  fusion/.style={lawsbox, fill=lawsNeutral, line width=0.9pt},
  decision/.style={lawsbox, fill=lawsAccent, line width=0.9pt},
%
  % --- archi ---
  arr/.style={
    -{Stealth[length=2mm, width=1.4mm]},
    draw=lawsInkSoft,
    line width=0.6pt,
  },
  busline/.style={draw=lawsInkSoft, line width=0.6pt},
%
  % --- etichette degli archi ---
  arrlabel/.style={
    font=\scriptsize,
    text=lawsInk,
    inner sep=1.5pt,
    fill=white,
  },
  % --- annotazione laterale (statuto epistemico, note di lettura) ---
  sidenote/.style={
    font=\scriptsize\itshape,
    text=lawsInkSoft,
    align=left,
    inner sep=1pt,
  },
}
%